% !TEX root = ../thesis.tex

\chapter{Conclusion and Future Work}
\label{ch:conclusion}

% Expansion of proposal §7 (Conclusion). Three top-level sections:
% Contributions, Direct RQ Answers, Future Work.

\section{Summary of Contributions}
\label{sec:conclusion:contributions}

% Migrated and lightly edited from proposal §7 with thesis-specific scope
% (four generators rather than three).
This thesis delivers a systematic and rigorously controlled comparison of
GAN-based and LLM-based methods for privacy-preserving synthetic medical data
generation. By evaluating HealthGAN, PATE-GAN, Pythia-70M, Pythia-1B, and
Pythia-1B-DP under a unified experimental framework on the
\emph{Diabetes 130-US Hospitals} dataset~\cite{diabetes_db}, the study
establishes a transparent basis for analysing the privacy-utility trade-off
across fundamentally different generative paradigms. The integration of aligned 
training conditions, consistent preprocessing, and a twelve-metric evaluation 
pipeline ensures comparability across models. It helps ensure that differences 
in performance reflect model characteristics rather than experimental variability.

The four concrete contributions of this work are the following.

\paragraph{Unified evaluation pipeline:} A reproducible two-pillar framework 
is used. It comprises \emph{utility}, split into statistical similarity to 
the real data and predictive performance of downstream models, and \emph{privacy}. 
The framework is instantiated as twelve single-number metrics applied identically to
 every generator under comparison.

\paragraph{Head-to-head benchmark:} Direct comparison of five generators
spanning the GAN/LLM, DP/non-DP, and LLM-scale axes:
HealthGAN~\cite{healthgan}, PATE-GAN~\cite{pategan}, Pythia-70M-LoRA and
Pythia-1B-LoRA~\cite{Biderman2023Pythia,Hu2022LoRA,Miletic2024}, and
Pythia-1B-DP-LoRA (Opacus DP-SGD~\cite{Yousefpour2021Opacus,Abadi2016DPSGD}).
The Pythia-1B non-DP scale anchor and the Pythia-1B-DP variant together 
constitute a scope extension beyond the originally submitted proposal. 
They yield the first within-family DP-cost comparison for tabular LLM 
synthesis on this benchmark at an architecturally controlled scale.

\paragraph{Within-family DP cost characterisation:} Quantitative evidence 
is provided on the marginal cost of formal $(\varepsilon = 5.0, 
\delta = 10^{-5})$-differential privacy. This cost is measured 
separately for the GAN family, from HealthGAN to PATE-GAN, and 
for the LLM family, from Pythia-1B to Pythia-1B-DP. On the LLM 
side, the toggle is held at a fixed model scale to isolate the 
DP-SGD effect from a confounding scale change.

\paragraph{Practical guidance:} A decision-oriented summary identifies when 
a non-DP GAN is sufficient. It also indicates which differentially 
private mechanism to prefer when a formal guarantee is required. 
Finally, it shows where LLM-based synthesis remains competitive 
under privacy budgets and governance constraints.

\section{Answers to the Research Questions}
\label{sec:conclusion:rq-answers}

% Condensed restatement of the "Answer to RQx" paragraphs from
% §\ref{sec:results:rq1:answer}, §\ref{sec:results:rq2:answer}, and
% §\ref{sec:results:rq3:answer}; the full supporting analysis stays in
% Chapter~\ref{ch:results}.
The three research questions stated in Section~\ref{sec:intro:rqs} are
answered directly below; the complete numerical evidence and figures
that support each answer are reported in Chapter~\ref{ch:results}.

For the first research question, the evaluation problem is answered by 
the two-pillar framework operationalised in Section~\ref{sec:method:evaluation}. 
Utility is assessed through statistical similarity and downstream 
predictive performance. Privacy is assessed separately through 
instance-level distance metrics and membership- and attribute-inference 
simulations. Instantiated as twelve single-number metrics (Table~\ref{tab:rq1-coverage}), 
the framework captures distinct evaluation dimensions. These branches 
are jointly necessary rather than individually sufficient. This is 
illustrated by PATE-GAN's degenerate distance scores alongside its 
utility collapse. A recommended minimal subset covering both pillars without
redundancy is given in Section~\ref{sec:results:rq1:minimal-suite}.

For the second research question, the privacy-utility comparison shows that
neither family dominates outright, but each occupies a different region of the
trade-off plane. Among the guarantee-free generators, the choice is a genuine 
trade-off. HealthGAN gives the best marginal fidelity and the only distance 
ratio at the NNDR unity reference. In contrast, the LLM generators, 
especially Pythia-1B and then Pythia-70M, give the highest statistical 
and predictive utility. Their performance is closest to the real-data ceiling, 
with an AUC of $0.682$ against $0.701$. Once $(\varepsilon = 5.0,
\delta = 10^{-5})$-differential privacy is demanded, the balance
tips decisively toward the LLM family. Pythia-1B-DP remains usable
while removing its non-private counterpart's observed attribute-inference
leak. In contrast, the differentially private GAN as implemented here,
PATE-GAN, collapses on both utility branches and earns only the apparent
privacy of off-manifold output (Section~\ref{sec:results:rq2:answer}).

For the third research question, the matched-budget comparison shows that the
cost of enforcing differential privacy at
$(\varepsilon = 5.0, \delta = 10^{-5})$ is strongly family-dependent. 
The LLM-family toggle from Pythia-1B to Pythia-1B-DP provides a clean test 
on a shared backbone. It costs only $0.0991$ in TSTR AUC-ROC while leaving 
statistical similarity and survival agreement essentially intact. It also 
converts that moderate cost into a genuine privacy gain by removing the
observed attribute-inference leak and driving outlier linkability to zero. The
GAN-family toggle (HealthGAN $\to$ PATE-GAN) instead collapses the utility
pillar for no genuine privacy gain, and must be read with the caveat that it
conflates the DP mechanism with an architecture change rather than isolating
differential privacy
(Section~\ref{sec:discussion:limit:gan-toggle-asymmetry}). The conditional
recommendation is therefore clear: when a formal guarantee at $\varepsilon =
5.0$ is required on tabular medical data of this kind, differentially private
LLM fine-tuning through DP-SGD yields a generator that remains usable, whereas
the differentially private GAN considered here does not.

\section{Future Work}
\label{sec:conclusion:future}

% Concrete extensions the thesis does not pursue but flags as natural
% follow-ups, each tied to a finding or limitation of the present study.
The findings reported above suggest several concrete extensions that this
thesis does not pursue but that follow naturally from its scope and
limitations.

\subsection{Larger Language Models}
\label{sec:conclusion:future:llm}

Pythia-1B attains the highest utility of any generator yet still trails the
real-data ceiling, and DP-SGD failed to converge usefully at the Pythia-70M
scale (Section~\ref{sec:results:rq3:pythia-70m-failure}). Both observations
motivate fine-tuning larger autoregressive models, such as Pythia-1.4B or
Llama-3-8B, under the same low-rank adaptation used here. Synthesis quality
improves with model size~\cite{Miletic2024} and larger models absorb DP-SGD
noise more gracefully~\cite{Li2022LargeDP}. So a larger backbone may both
narrow the residual utility gap and reduce the differential-privacy cost, at
the price of the compute and memory overhead of per-sample-gradient DP-SGD at
that scale.

\subsection{Cross-Dataset Generalisation}
\label{sec:conclusion:future:cross-data}

All results are obtained on a single benchmark~\cite{diabetes_db}, so the
generator rankings cannot be assumed to transfer to other clinical settings.
Repeating the unified pipeline on richer EHR corpora such as MIMIC-IV~\cite{mimic} 
or eICU would test the generality of these findings. These corpora contain 
more continuous variables, higher attribute cardinality, and explicit temporal 
structure. Such experiments would establish whether the LLM family retains its 
utility advantage and whether the GAN-side differential-privacy collapse is 
dataset-specific or systematic.

\subsection{Privacy-Budget Ablation}
\label{sec:conclusion:future:epsilon}

The privacy comparison is conducted at a 
single matched budget of $(\varepsilon = 5.0, 
\delta = 10^{-5})$. This gives one point on each 
family's privacy-utility curve rather than its slope. Repeating PATE-GAN and 
Pythia-1B-DP across a range of budgets, such as $\varepsilon \in \{1, 3, 5, 10\}$~\cite{Nahid2024},
would reveal whether the GAN-family collapse persists at looser budgets. 
It would also locate the budget at which the LLM family's utility cost becomes prohibitive.

\subsection{Stronger Privacy Adversaries}
\label{sec:conclusion:future:adversaries}

The privacy pillar relies on a distance-based membership-inference adversary
and a predictor-based attribute-inference attack. Both are comparatively weak
threat models, so the negligible leakage they report
(Section~\ref{sec:results:rq2:privacy}) is a lower bound on disclosure risk
rather than a guarantee. Future work should evaluate stronger adversaries, 
including shadow-model membership inference~\cite{Shokri2017MIA}, training-data 
extraction~\cite{Carlini2021Extracting}, and white-box attacks. These attacks 
would test that bound. By comparing the differentially private variants against 
their non-private counterparts, such work could show whether the formal guarantee 
delivers protection that the present attacks are too weak to detect.

\subsection{Longitudinal and Multimodal Records}
\label{sec:conclusion:future:longitudinal}
This study treats each encounter as one independent cross-sectional row. 
Extending the comparison to time-series synthesis~\cite{dash2020medical,bhanot2022investigating} 
and to multimodal records combining tabular fields with imaging or 
free-text notes would test the generality of the findings. It would show whether the privacy-utility 
conclusions hold once temporal and cross-modal dependencies must also be preserved.

\subsection{Non-Linear Dependency Metrics}
\label{sec:conclusion:future:nonlinear}

The statistical-similarity branch quantifies joint structure through the linear
correlation difference. Extending it with a non-linear measure such as a
mutual-information matrix difference~\cite{kaabachi2025scoping} would capture
dependencies that linear correlation cannot.

\section{Closing Remarks}
\label{sec:conclusion:closing}

% Migrated and lightly adjusted from proposal §7 closing paragraphs.
This thesis combines utility assessment with instance-level privacy assessment. 
Utility spans both distributional similarity to the real data and predictive 
performance of downstream models. This combined evaluation determines not only 
which models perform best, but also under which conditions each class of models 
is most appropriate. This enables a nuanced interpretation of when non-DP 
synthesis is sufficient. It also clarifies when formal differential privacy 
is necessary. Finally, it shows how LLM-based generators behave in terms of 
memorisation and disclosure risk in clinical settings. Beyond generating 
comparative insights, the thesis contributes a
reusable, methodologically grounded evaluation protocol that can support
future research and practical governance. Although the study focuses on
tabular EHR-like data and a defined set of generative models, the framework
provides a foundation for subsequent extensions to longitudinal records,
multimodal data, and emerging DP-enhanced LLMs. Ultimately, the findings are
intended to inform model selection, guide responsible data-release strategies,
and strengthen the methodological basis for safe and high-utility synthetic
data use in healthcare research.

\cleardoublepage
